\section{Experimental Design}
\label{app:experimental-design}

Our experiments explore the hyperparameter loss surface across scales via random search.

\paragraph{Models.} We measure model size by the \textit{effective parameter count}, parameters adjusted for FLOPs, so that compute is $c = 6pd$ for $p$ effective parameters and $d$ tokens. This count excludes the embedding layer (which uses no FLOPs) but includes the unembedding layer and the cost of attention \citep[\S B]{Porian2024Resolving}. We target parameter counts at increasing powers of two: $2^{22}\approx4\text{M}$ to $2^{28}\approx268\text{M}$, and construct models at each scale with two experimental designs: a hand-crafted model ladder and randomly sampled architectures.

\paragraph{Training and evaluation.} We train with a warmup-stable-decay (WSD) schedule \citep{hu2024minicpm}. Following \citet{Hagele2024ScalingLaws}, we reuse the stable phase from the main run and branch off decay phases from checkpoints taken $\sfrac{1}{8}, \sfrac{2}{8}, \ldots, \sfrac{8}{8}$ through training. We evaluate checkpoints before and after decaying the learning rate, thus a single sweep yields stable and decayed evaluations across eight token budgets at a fraction of the cost.

\paragraph{Analysis.} We partition the scales into training (4M--34M), validation (67M--134M), and test (268M), so reported numbers measure extrapolation from small models to larger held-out ones. By default, we fit the joint scaling law of Equation~\ref{eq:scaling-law_joint} to the best loss attained at each parameter--data budget pair. Like prior work \citep{hilton2023scalinglawssingleagentreinforcement, li2025misfitting}, we found shorter training runs fit the scaling law poorly, therefore we discard the two smallest token budgets which noticeably improves performance on our validation set.


\subsection{Modeling}
\label{app:experimental-design:modeling}

The modeling recipe we study is a Llama variant, as implemented in Meta Lingua: \url{https://github.com/facebookresearch/lingua} (commit: 437d680) \citep{meta-lingua}. Hyperparameters not specified here were either left to Lingua's defaults or swept over as described in \Cref{app:experimental-design:search-distributions}.

\paragraph{Data.} We pretrain on the 100B token subset of FineWeb-Edu \citep{Penedo2024TheFineWeb}. We use the subset to save storage and preprocessing since the scales we consider require less than 100B tokens. We split this subset into training and validation sets using Lingua's built-in preprocessing pipeline with 48GB of memory and Lingua's suggested seed of 42. We create two such datasets: one shuffling all the data in 1 chunk for single GPU training (4M--134M parameters), and another splitting the data into 8 chunks for 8 GPU training (268M parameters).

\paragraph{Architecture.} For the tokenizer, we use \texttt{p50k\_base} from \texttt{tiktoken} (\url{https://github.com/openai/tiktoken}), the BPE tokenizer from GPT-3.5 (\texttt{text-davinci-003}), because it has an appropriately sized vocabulary for our scale. The context length is 1{,}024 with special tokens to distinguish the beginning and ending of each sequence. For the neural network, we use a decoder-only Llama model \citep{grattafiori2024llama3herdmodels} available in Lingua (\texttt{apps/main/train.py}) \citep{meta-lingua}. The model uses RMSNorm for normalization \citep{Zhang2019Root}, SwiGLU for the activation function \citep{shazeer2020gluvariantsimprovetransformer}, and RoPE for position embeddings \citep{Su2024RoFormer}. All runs use a pre-normalization architecture except for those explicitly specified as post-normalization (e.g., in \S\ref{sec:designing-small-scale-experiments}), which we implemented as described by \citet{Xiong2020OnLayer}. Aside from Lingua's defaults, we padded the vocabulary size (\texttt{model.vocab\_size}) to 50{,}536. In addition, we set \texttt{model.multiple\_of} to 8 since the default value of 256 is larger than our smaller models' hidden dimensions; 8 compromises between allowing a range of hidden dimensions at small scales while partially exploiting increased GPU efficiency.

\paragraph{Optimization.} For optimization, we use AdamW \citep{loshchilov2018decoupled} with a warmup-stable-decay (WSD) learning rate schedule \citep{hu2024minicpm}. Unlike Lingua, we use a linear cooldown for the WSD schedule decaying to a \texttt{optim.lr\_min\_ratio} of 1e{-}6 times the peak learning rate. Each scale we train is defined by a target number of effective parameters (described in \Cref{app:experimental-design:search-distributions}), ranging from $2^{22} \approx 4M$ to $2^{28} \approx 268M$ in powers of two. Training continues until the number of tokens equals 32 times the target number of effective parameters. The total tokens divided by the tokens per gradient step determines the number of steps. The effective parameter count is essentially the same as FLOPs per step as used in \citet{Hoffmann2022AnEmpirical} (Approaches 1 \& 2), thus we train about 50\% further than the Chinchilla-optimal ratio of $\sim$20 tokens-per-parameter, allowing us to explore both the under- and over-trained regimes. Following \citet{Hagele2024ScalingLaws}, we first train each configuration on a main run without decaying the learning rate, saving checkpoints $\sfrac{1}{8}, \sfrac{2}{8}, \ldots, \sfrac{8}{8}$ through training; afterwards we decay each of those checkpoints by extending training with 25\% more steps, thus the final results have their learning rates decayed over the last 20\% of training. The initial checkpoints provide our constant learning rate results (\emph{stable}), whereas the decayed checkpoints provide our results for the full schedule (\emph{decay}); all results include learning rate warmup.

\paragraph{Evaluation.} For pretraining, we evaluate the model via bits-per-character (BPC), which normalizes the pretraining loss for differences in tokenization. To compute BPC, we change the logarithm's base on the negative log-likelihood per character already provided by Lingua. We calculate BPC on the validation set, limiting \texttt{eval.validation.max\_steps} to 1{,}000. For capabilities, we evaluate the model using Lingua's integration with the EleutherAI \texttt{lm-evaluation-harness} (v0.4.11) \citep{eval-harness}. We evaluate on six multiple-choice tasks: AI2 ARC (Easy) and (Challenge) \citep{clark2018thinksolvedquestionanswering}, BoolQ \citep{clark-etal-2019-boolq}, MMLU \citep{hendrycks2021measuring}, OpenBookQA \citep{mihaylov-etal-2018-suit}, and PIQA \citep{Bisk2020PIQA}. Our evaluations limit \texttt{eval.generator.max\_tokens} to 16{,}384 but otherwise use Lingua's defaults.

\paragraph{Random Seeds.} We fix the following random seeds made available by Lingua: \texttt{seed} to 0, \texttt{data\-.seed} to 1, \texttt{model\-.seed} to 2, \texttt{eval\-.harness\-.random\_seed} to 3, \texttt{eval\-.harness\-.numpy\_random\_seed} to 4, \texttt{eval\-.harness\-.torch\_random\_seed} to 5, \texttt{eval\-.harness\-.fewshot\_random\_seed} to 6. Accordingly, data order and other details were kept constant while runs differed primarily in their hyperparameters.

\subsection{Search Distributions}
\label{app:experimental-design:search-distributions}

Our results analyze three different experiments, each defined by a search distribution over the hyperparameters.

All the experiments randomly sample configurations across a range of scales, with each scale defined by the \emph{effective} number of parameters. The effective number of parameters is proportional to the FLOPs per token so that $c = 6pd$ when $c$ is compute, $p$ the effective parameters, and $d$ the number of tokens. Following \citet[\S B]{Porian2024Resolving}, we compute it as:
%
\begin{equation}\label{eq:effective-number-of-parameters}
    N_{eff} \coloneqq \left(3 d_{ff} + 4d\right)dl + dv + ndl
\end{equation}
%
where $d_{ff}$ is the feedforward dimension, $d$ the hidden dimension, $l$ the number of layers, $v$ the vocabulary size, and $n$ the context length of the transformer. In particular, parameters in the embedding layer are not counted because they do not contribute FLOPs, while parameters in the unembedding layer are counted because they do. See \citet[\S B]{Porian2024Resolving} for an overview of different ways to count parameters.

\begin{figure}
    \centering
    \begin{align*}
        \hparam{batch\_size}   &\sim \distribution{DiscreteUniform}\left(\{64, 128, 256, \ldots, 4{,}096\}\right) \\
        \hparam{lr}            &\sim \distribution{LogUniform}(1e{-}5, 1e{-1}) \\
        \hparam{beta1}         &\sim \distribution{LogitUniform}(0.7, 0.999) \\
        \hparam{beta2}         &\sim \distribution{LogitUniform}(0.9, 0.9999) \\
        \hparam{warmup}        &= \operatorname{round}\left(\hparam{steps} \times P\right),\quad P \sim \distribution{LogUniform}\left(1e{-}3, \sfrac{1}{8}\right) \\
        \hparam{weight\_decay} &\sim \distribution{LogUniform}(1e{-}4, 1e0) \\
        \hparam{rope\_theta}   &\sim \distribution{LogUniform}\left(2^{10}, 2^{20}\right) \\
    \end{align*}
    \caption{
        \emph{The search distribution for all non-architectural hyperparameters.} The log uniform and logit uniform distributions are obtained by taking the log or logit of the bounds, sampling uniformly on this scale, then inverting the transformation back to the original scale.
    }
    \label{fig:search-distribution_non-architectural}
\end{figure}

To isolate the effect of tuning the architecture, configurations of the same index share all non-architectural hyperparameters across the different designs. Thus, all differences between the experiments come from the architecture alone. The search distribution over the shared hyperparameters is given in \Cref{fig:search-distribution_non-architectural}. Typically, researchers search hyperparameters like the learning rate or weight decay on a log scale; using the proper scale for each hyperparameter is also important in efficiently reaching the noisy quadratic limit \citep{lourie2025hyperparameter}. In practice, we found the simple heuristic of sampling real values on a linear scale, positive values on a log scale, and bounded values on a logit scale worked quite well. The warm up is sampled as a proportion of the total training steps up to the first $\sfrac{1}{8}$ of training so that the first checkpoint (and evaluation) is always fully warmed up.

\begin{table}
    \centering
    \begin{tabular}{c|rrrr}
        \toprule
          $N_{eff}$ & $\hparam{dim}$ & $\hparam{n\_layers}$ & $\hparam{head\_dim}$ & $\hparam{n\_heads}$ \\
        \midrule
              3.47M &             64 &                    2 &                   64 &                   1 \\
              7.78M &            128 &                    4 &                   64 &                   2 \\
              17.1M &            256 &                    4 &                   64 &                   4 \\
              32.4M &            384 &                    6 &                   64 &                   6 \\
              62.6M &            512 &                   10 &                   64 &                   8 \\
               133M &            768 &                   12 &                   64 &                  12 \\
               270M &          1,024 &                   16 &                   64 &                  16 \\
        \bottomrule
    \end{tabular}
    \caption{
        \emph{The model ladder used for architectural hyperparameters in the \texttt{ladder} experiment.} The leftmost column gives the actual effective number of parameters which differs slightly from the targeted number due to architectural constraints.
    }
    \label{tab:model-ladder}
\end{table}

\paragraph{Model Ladder (\texttt{ladder}).} We sample configurations for target scales starting from $N_{eff} = 2^{22} \approx 4M$ up to $2^{28} \approx 268M$ in powers of two. Excluding 2 failed runs, we search 867 configurations at 4M effective parameters, 826 at 8M, 867 at 17M, 934 at 34M, 128 at 67M, 128 at 134M, and 64 at 268M. For the \texttt{ladder} experiment, we follow the common practice of setting architectural hyperparameters via a hand-crafted model ladder \citep{Hagele2024ScalingLaws, bhagia2025establishing}. The hidden dimension, number of layers, and number of attention heads at each scale were set according to Table~\ref{tab:model-ladder}. We increase the width and depth of successive models in order to steadily increase their parameter count, while keeping other architectural elements (such as the $\hparam{head\_dim}$) constant. Following \citet{Porian2024Resolving}, we keep the aspect ratio between 32 and 64. The transformer architecture is subject to a number of constraints, depending on the implementation (e.g., $\hparam{head\_dim}$ divides $\hparam{dim}$, $\hparam{dim}$ is used both for the hidden and the embedding dimension, matrix dimensions should be multiples of 8 when running on a GPU, and so on). These constraints must be considered when designing a ladder.

\begin{figure}
    \centering
    \begin{align*}
        \hparam{n\_layers} &= \operatorname{round}\left(U^2\right),\quad U \sim\distribution{Uniform}\left(\sqrt{2}, \sqrt{48}\right) \\
        \hparam{n\_heads}  &\sim \distribution{DiscreteProportional}\left(\left\{k \mid k \text{ divides } \sfrac{\hparam{dim}}{8},\; 1 \leq k < \min\left(24, \sfrac{\hparam{dim}}{8}\right)\right\}\right) \\
    \end{align*}
    \caption{
        \emph{The search distribution for architectural hyperparameters used in the \texttt{random}, \texttt{prenorm}, and \texttt{postnorm} experiments.} The discrete proportional distribution samples from a finite set of positive numbers with probability proportional to each number. By construction, $\hparam{dim}$ is always divisible by $\text{\texttt{model.multiple\_of}} = 8$.
    }
    \label{fig:search-distribution_architectural}
\end{figure}

\paragraph{Random Architecture (\texttt{random}).} We sample configurations for target scales starting from $N_{eff} = 2^{22} \approx 4M$ up to $2^{27} \approx 134M$ in powers of two. Excluding 2 failed runs, we search 128 configurations at 4M effective parameters, 126 at 8M, 128 at 17M, 128 at 34M, 128 at 67M, and 128 at 134M. In the \texttt{random} experiment, we also tune the architectural hyperparameters via random search. The number of layers and attention heads for each configuration are sampled from the distribution described in \Cref{fig:search-distribution_architectural}. Given the effective parameters, the number of layers determines the hidden dimension via \Cref{eq:effective-number-of-parameters}, which is a quadratic in $d$ given the other parameters and the fact that $d_{ff} = \sfrac{8}{3} d$ for our model. We round the solution to the nearest multiple of 8 for GPU efficiency. Once we have the hidden dimension, the number of heads determines the head dimension: $\hparam{dim} = \hparam{head\_dim} \times \hparam{n\_heads}$. We ensure $\hparam{n\_layers}$ is at least 2 because induction heads cannot form with fewer layers \citep{olsson2022incontextlearninginductionheads}; we also use a square root scale because we found it balanced making both the number of layers and the hidden dimension more uniform. $\hparam{n\_heads}$ must divide the hidden dimension, so we sample proper divisors of $\sfrac{\hparam{dim}}{8}$ to guarantee that $\hparam{head\_dim}$ is even (required by RoPE) and greater than 16 (required by the attention implementation). Since small divisors are much more common, we bias towards larger ones by sampling divisors with probability proportional to their size. As previously mentioned, transformers are subject to a number of constraints that must be accounted for when sampling architectural hyperparameters: $\hparam{head\_dim} \geq 16$ for common attention implementations, $\hparam{head\_dim}$ must be even for RoPE, $\hparam{head\_dim}$ divides $\hparam{dim}$, $\hparam{dim}$ provides both the hidden and the embedding dimension, and matrix dimensions should be multiples of 8 for GPU efficiency. These constraints complicate designing a good search distribution;\footnote{
    A typical model ladder avoids most of these complications, as they primarily appear as edge cases in small networks.
} the distribution in \Cref{fig:search-distribution_architectural} has been constructed to obey them.

\paragraph{Pre-Normalization (\texttt{prenorm}) vs. Post-Normalization (\texttt{postnorm}).} We sample configurations for target scales of $N_{eff} = 2^{22} \approx 4M$, $2^{25} \approx 34M$, and $2^{27} \approx 134M$. For \texttt{prenorm}, we reuse all the results at those scales from the \texttt{random} experiment. For \texttt{postnorm}, we reuse the initial configurations of the \texttt{random} experiment (including architectural parameters), but we re-run them after changing the transformer to a post-normalization architecture. In addition, we expand the number of configurations due to the post-normalization architecture's increased tuning difficulty. Excluding 1 failed run, this expansion results in 511 configurations at 4M effective parameters, 512 at 34M, and 128 at 134M.


\subsection{Analysis}
\label{app:experimental-design:analysis}

\paragraph{Noisy Quadratic Limit.} We generate confidence bands for the eCDF using the LD Highest Density bands from \citet{lourie-etal-2024-show}, and we fit the noisy quadratic distribution as described in \citet{lourie2025hyperparameter}, both using \texttt{opda} (v0.8.0): \url{https://github.com/nicholaslourie/opda}. To fit the noisy quadratic limit, we set the threshold $\theta$ at various percentiles of the score distribution. We used different percentiles for each parameter--data budget, with the size of the asymptotic regime generally increasing in both. In the \texttt{ladder} and \texttt{random} / \texttt{prenorm}) experiments, they were as follows: for 4M, 8M, and 17M effective parameters, 7\% at checkpoints 1--6 and 8\% at checkpoints 7--8; for 34M, 7\% for checkpoints 1--2, 8\% at 3, 10\% at 4--6, and 12\% at 7--8; for 67M, 12\% at 1, 15\% at 2, 20\% at 3, 25\% at 4, 30\% at 5, 32\% at 6, and 35\% at 7--8; for 134M and 268M, 12\% at 1, 15\% at 2, 20\% at 3, 25\% at 4, 30\% at 5, 35\% at 6, 40\% at 7, and 50\% at 8. In the \texttt{postnorm} experiment, they were as follows: for 4M effective parameters, 3\% at checkpoint 1, 4\% at 2, 5\% at 3, 6\% at 4, and 7\% at 5--8; for 34M and 134M, 7\% at 1--2, 8\% at 3, and 10\% at 4--8. We restricted $\alpha$ to be positive (as bits-per-character always is), and we restricted $\gamma$ to be 1 to at most the nominal number of hyperparameters (7 for \texttt{ladder} and 9 for \texttt{random} / \texttt{prenorm} and \texttt{postnorm}).

\paragraph{Scaling Laws.} To fit and evaluate the scaling laws, we partition the scales into training (4M--34M), validation (67M--134M), and testing (268M); thus, we evaluate the scaling laws based on their ability to extrapolate from small to large models. We used training to fit the laws, validation to make decisions during our research, and testing for the final evaluation. For the scaling law, we fit \Cref{eq:scaling-law_joint} to the best validation loss from each parameter--data budget pair, or (where explicitly stated) to the validation losses from the run achieving the minimum final loss at each scale. In both cases, we drop the first two checkpoints ($\sfrac{1}{8}$ and $\sfrac{2}{8}$ of training) as it substantially improved the fit \citep{hilton2023scalinglawssingleagentreinforcement}. Where specified, we also tie the scaling exponents together ($\iota = \kappa$). To perform the optimization, we minimized the mean squared error via SciPy's \texttt{scipy.optimize.differential\_evolution} \citep{SciPy2020Virtanen}, with bounds of $[0, 1]$ for $\epsilon$, $[0, 10]$ for $\ln \zeta$ and $\ln \eta$ (which we fit on log scales), and $[0, 1]$ for $\iota$ and $\kappa$. We used a \texttt{popsize} of 120 except when fitting laws jointly with a single irreducible error for which we used a \texttt{popsize} of 60 due to the higher dimensionality and longer runtime. To evaluate the scaling laws, we report mean squared error between the predicted and observed loss on the held-out scales, with held-out scales using the same preprocessing as those in training (e.g., whether the best run or the best configuration for each parameter--data budget is used).

\paragraph{Interpretable Models of the Hyperparameter Loss Surface.} In \Cref{app:interpretable-models-of-the-hyperparameter-loss-surface}, we compare various models of the hyperparameter loss surface. These models are fit to losses generated via the same methodology used for scaling laws (best loss for each parameter--data budget, dropping the first two checkpoints), except instead of selecting the loss for the single best configuration, we take all configurations within the asymptotic regime as defined by the fractions used for our noisy quadratic limits. Before fitting the models, we transform the hyperparameters to put them on an appropriate scale: a log-scale for $\hparam{n\_effective\_parameters}$, $\hparam{n\_tokens}$, $\hparam{batch\_size}$, $\hparam{lr}$, $\hparam{weight\_decay}$, $\hparam{dim}$, $\hparam{n\_layers}$, $\hparam{n\_heads}$, and $\hparam{rope\_theta}$; a log-plus-one-scale for $\hparam{warmup}$ (to avoid log of zero); and a logit-scale for $\hparam{beta1}$ and $\hparam{beta2}$. The interpretable models are fit by first taking a good initialization and then improving it with SciPy's \texttt{scipy.optimize.minimize} using the \texttt{"Nelder-Mead"}, \texttt{"Powell"}, and \texttt{"BFGS"} methods in sequence. For the initialization, we fit a quadratic polynomial via ordinary least squares, algebraically solve for the equivalent estimate in the format of \Cref{eq:static-empirical-model_projection}, and then project the Hessian to a lower rank ($\gamma = 3$) by dropping the eigenvalues of least absolute value. We then use the \texttt{static} model to initialize the \texttt{dynamic optima} and the \texttt{dynamic sensitivity} models, setting the new components to zero since the models are nested. We similarly initialize the \texttt{dynamic} model from the \texttt{dynamic optima} model. For the \texttt{multilayer perceptron} model, we implement the neural network using scikit-learn \citep{scikit-learn}. The pipeline applies a quadratic basis expansion to the hyperparameters, normalizes the features via \texttt{preprocessing.RobustScaler}, and finally predicts via a multilayer perceptron (MLP) with \texttt{hidden\_layer\_sizes=[128]}, \texttt{learning\_rate="adaptive"}, \texttt{batch\_size=256}, and \texttt{max\_iter=4096}. We grid search \texttt{alpha} over 1e{-}4, 1e{-}3, 1e{-}2, 1e{-}1, 1e0 and \texttt{learning\_rate\_init} over 1e{-}5, 3e{-}5, 1e{-}4, 3e{-}4, 1e{-}3 using the train-validation split described previously. For the MLP, we fit the scaling law beforehand and then fit the MLP to the percent excess loss, rather than fitting them jointly.

\subsection{Implementation}
\label{app:experimental-design:implementation}

Our implementation further involves the following details.\footnote{
    Preliminary experiments were run on the NYU Torch cluster. We later expanded our random searches using Meta's compute infrastructure. This section describes the compute setup used at NYU, which is sufficient to reproduce our findings.
}

\paragraph{Training Pipeline.} We use gradient accumulation so that we can tune the batch size despite different scales requiring different amounts of memory and GPU setups (1 vs. 8 GPUs). We set the per-GPU batch size to the largest power of two that safely avoids out-of-memory errors in order to maximize throughput while also evenly dividing the number of tokens (16 sequences for 4M--67M parameters, 8 sequences for 134M--268M parameters). Each scale is assigned a base number of gradient accumulation steps chosen so that after accounting for the per-GPU batch size and the number of GPUs the model will process 64 sequences per update. We then multiply that base by a power of two to achieve the distribution described in \Cref{fig:search-distribution_non-architectural} for the total batch size. As a result, the batch size sweep is hardware-agnostic and always evenly divides the total number of training tokens. In addition, we set \texttt{gc\_collect\_freq} to 1{,}024, \texttt{data.prefetch\_size} to 1{,}024, \texttt{distributed.compile} to \texttt{true}, and \texttt{distributed.fsdp\_type} to \texttt{"full\_shard"} in order to make training efficient. Finally, we set \texttt{logging.freq} to 1 and found it necessary to set \texttt{env.ENABLE\_INTRA\_NODE\_COMM} to \texttt{"0"} to avoid an unresolved error in the PyTorch version.

\paragraph{Software.} We train the models using Meta Lingua (commit: 437d680), as previously described \citep{meta-lingua}. We modify the code base to support the 100B token subset of FineWeb-Edu, use a linear decay in the WSD schedule, and use PyTorch's Scaled Dot Product Attention (SDPA) instead of Flex Attention as Flex Attention would not successfully generate kernels for some of the transformer shapes that we randomly sample. In addition, we switch the transformer architecture from pre-norm to post-norm for some experiments as previously discussed. We isolate the software environment in an Apptainer (Singularity) container \citep{singularity2021}. The container runs Ubuntu 24.04.2 LTS with CUDA 12.8 (V12.8.93) and cuDNN 9.8, and we run Lingua via a Python 3.11.14 conda environment with PyTorch 2.7.0+cu128 \citep{Paszke2019PyTorch}.

\paragraph{Hardware.} Training was done on 80GB NVIDIA A100 GPUs on a SLURM cluster. Model sizes from 4M--134M effective parameters were trained with 1 GPU, while models with 268M effective parameters were trained data-parallel with 8 GPUs.

